\chapter{Input, output and utilities}\label{ch:io-utilities}
\chapterstatus{written}{2026-11-15}

Snapshot readers and writers, plain-text tables, parameter files, physical units, timing and memory instrumentation,
the performance suite, the Python interfaces and a few smaller utilities.

\begin{modulebox}{I/O}
\nsp{FML::FILEUTILS}: \file{GadgetUtils/GadgetUtils.h} (\cls{GadgetReader}, \cls{GadgetWriter}),
\file{GadgetUtils/GadgetUtilsHDF5.h}, \file{RamsesUtils/RamsesUtils.h} (\cls{RamsesReader}), \file{FileUtils/FileUtils.h}
(\code{loadtxt}, \ldots), \file{FileUtils/DataPath.h} (\code{fml_file}); the program \file{RamsesUtils/ramses2pk}.\par
Tests: \file{tests/test_gadget.cpp}, \file{tests/test_ramses.cpp}, \file{tests/test_utils.cpp}.
\end{modulebox}
\begin{modulebox}{Parameters and units}
\file{ParameterMap/ParameterMap.h} (\cls{ParameterMap}), \file{LuaFileParser/LuaFileParser.h} (\cls{LuaFileParser}),
\nsp{FML::UTILS}: \file{Units/Units.h} (\nsp{SI}, \cls{ConstantsAndUnits}).\par
Documentation: \file{Units/README.md}. Tests: \file{tests/test_lua.cpp}, \file{tests/test_utils.cpp}, \file{tests/test_units_odr.cpp}.
\end{modulebox}
\begin{modulebox}{Instrumentation and interfaces}
\file{Timing/Timer.h}, \file{Timing/Timings.h}, \file{MemoryLogging/MemoryLogging.h}, \file{Performance/} (counters and the
benchmark runner), \file{Plotting/Matplotlibcpp.h}, \file{Chunkpix/Chunkpix.h}, \file{PythonAnalysis/} (\code{fml_analysis}),
\file{MakePythonWrapper/}, \file{HaloModel/PythonWrapper/}, \file{Visualization/} (projections, colour maps, a PNG writer),
\file{ExampleData/} (small inputs generated in code).\par
Documentation: \file{Performance/README.md}, \file{PythonAnalysis/README.md}, \file{Visualization/README.md}.
\end{modulebox}

\section{Snapshot formats}\label{sec:io-snapshots}

\paragraph*{GADGET binary} \cite{Springel2005}. \cls{GadgetReader} reads format 1 and format 2 (blocks preceded by a label block)
files of either endianness (detected and swapped), positions and velocities in float or double, 4- or 8-byte IDs (from the block
size) and the MASS block (types with zero \code{MassTable} entry, stored with \code{set_mass}); empty files are skipped and unknown
format-2 blocks ignored. Positions are returned in $[0,1)$ (scaled by the box, with a factor for files in kpc/$h$), velocities as
peculiar velocities in km/s (the file holds $v/\sqrt a$), the family/type is set if the particle has it (otherwise all are dark
matter), and with MPI each task can keep only the particles in its domain while reading all files. \cls{GadgetWriter} writes one file
per task with 4-byte IDs unless an ID or the total does not fit, several species per file with their \code{MassTable} masses
($\Omega_{\rm type}\rho_{\rm crit}B^3/N_{\rm type}$) or per-particle masses; blocks above $2^{31}$ bytes (the limit of the format) are an
error.

\paragraph*{GADGET HDF5.} \file{GadgetUtilsHDF5.h} (with \code{USE_HDF5}) reads and writes the HDF5 layout of GADGET-2/3/4, AREPO,
GIZMO and SWIFT (\code{/Header} attributes and \code{/PartTypeN/Coordinates}, \code{Velocities}, \code{ParticleIDs}, \code{Masses}), in
one or several files, with the same interface and conventions; velocity units in the file are converted to km/s.

\paragraph*{RAMSES} \cite{Teyssier2002}. \cls{RamsesReader} reads dark-matter particle files (positions, velocities in km/s from
\code{unit_l}/\code{unit_t}, masses, IDs, levels, families and tags as the particle provides them) from the header or the
\file{part_file_descriptor.txt}, keeping only the local domain if asked, and writes the \file{ic_deltab} initial-condition file.
\file{ramses2pk} computes power spectra directly from a RAMSES output.

\paragraph*{Plain text.} \code{loadtxt} and its variants read whitespace-separated numbers (comments with \#, ! or /; \code{nan},
\code{inf} and denormals), with a fixed or variable number of columns, a subsample fraction, or two extracted columns; malformed lines
are errors with the line number. \code{fml_file} resolves data files of the library from \code{FML_DIR}, the compiled
\code{FML_ROOT_DIR} or a relative path.

Grids and particle containers have raw checkpoints (\code{dump_to_file}, \code{load_from_file}; \cref{ch:grids-fft,ch:particles}).

\section{Parameters}\label{sec:io-parameters}

\cls{ParameterMap} is a map from names to typed values (string, int, bool, double and vectors of int and double, a
\code{std::variant}): \code{get<T>(name)} throws if the parameter is missing or of another type, and \code{get<T>(name, fiducial)}
returns the fiducial value only if the parameter is missing (a wrong type is still an error). \cls{LuaFileParser} reads Lua scripts as
parameter files (the format of \file{SimulationCode}, \cref{ch:simulation-code}): \code{read_int}, \code{read_double},
\code{read_string}, \code{read_bool} and arrays, each \code{required} or \code{optional} with a default. A parameter that is set with the
wrong type (1 for a boolean, 2.5 for an integer, a string in an array of numbers) is always an error, and reading many unset optional
parameters no longer corrupts the Lua stack. Being Lua, parameter files can compute values and use loops and conditions.

\section{Units and constants}\label{sec:io-units}

\file{Units.h} is the single source of physical constants: \nsp{FML::UTILS::SI} holds compile-time SI values (units, derived and
astronomical units, fundamental constants, recombination energies and rates, $H_0/h$ in s$^{-1}$; the ionisation energies are in joules),
and \cls{ConstantsAndUnits} gives all constants in a chosen system (\code{"SI"}, \code{"Planck"}, \code{"ParticlePhysics"},
\code{"Cosmology"}: Mpc, Hubble time, solar masses and kelvin, or user factors), derived from the same table. The Einstein--Boltzmann
solver's \code{Constants} namespace is an alias to it (\cref{ch:conventions}). The year is the 365.25-day Julian year.

\section{Instrumentation}\label{sec:io-instrumentation}

\paragraph*{Timing.} \cls{Timer} is a stopwatch; \cls{Timings} accumulates labelled timings over restarts (a repeated end does not count
twice) and prints a summary; \file{SimulationCode} reports the time of every part of a step with it.

\paragraph*{Memory.} With \code{MEMORY_LOGGING} the containers of the library use a logging allocator that records every allocation above a
minimum size with its label (grids and particle containers label themselves), keeps the current and peak memory and their history, and
prints the range over tasks; it is thread safe and bounded in the number of tracked allocations. \file{SystemMemory.h} gives the resident
set size (\cref{sec:par-communication}).

\paragraph*{Performance suite.} \file{Performance/} measures representative operations with full provenance, separately from the
correctness tests: an FFT round trip, grid copies and empty clones, Fourier resampling, particle migration and selected bispectrum
configurations, each validated outside the timed region. The runner (\file{run.py}) runs every case, rank and thread count in a fresh
process and stores JSON reports, raw logs, the argument vectors and a manifest (host, CPU affinity, OpenMP settings, the executable's
SHA-256). Opt-in counters (\code{FML_ENABLE_PERFORMANCE_COUNTERS}) count FFTW plan creations, transforms, grid allocations and copies, and
the payload bytes and messages of the transports.

\section{Python}\label{sec:io-python}

\begin{itemize}
\item \code{fml_analysis} (\file{PythonAnalysis/}, pybind11): a serial catalogue interface to the spatial statistics of
      \cref{ch:configuration-space}: \code{Catalogue(positions, boxsize)} with $k$-nearest-neighbour distances and CDFs, sphere counts and
      counts-in-cells summaries for query points, on a KD-tree, with NumPy arrays in and out, copies of all inputs and the GIL released during
      the C++ work. It does not initialise MPI or FFTW.
\item The halo-model wrapper (\file{HaloModel/PythonWrapper/}, pybind11) exposes \cls{HaloModel} (\cref{ch:halo-model}).
\item \file{MakePythonWrapper/} shows how to wrap any FML function with SWIG and NumPy arrays.
\end{itemize}

\section{Pictures}\label{sec:io-visualization}

\file{Visualization} makes pictures of fields and particles from C++ with no dependencies, for a quick look at a run or
frames of a movie. \code{project_grid} projects a slab of an \cls{FFTWGrid} or \cls{PencilGrid} along an axis (the sum or
the mean over the cells whose positions $i/N_{\rm mesh}$ are in a periodic slab) and \code{project_particles} the particles
(NGP or CIC, mass weighted when they have masses), giving the same \cls{Image} on every task. \cls{ColorScale} maps the
values linearly, logarithmically or symmetrically about zero, with fixed limits or percentiles of the image, onto a colour
map: magma, inferno, viridis, plasma and cividis (the matplotlib tables), coolwarm (diverging \cite{Moreland2009}), grey
and cubehelix \cite{Green2011} (a suffix \code{_r} reverses them). \code{write_png} has its own deflate (per-row filters,
LZ77 and fixed Huffman codes, stored blocks for noise), and \code{write_npy} writes the values for numpy.
\file{Analysis/Images.h} builds on it for the images of \file{SimulationCode} and SnapshotAnalyzer
(\cref{sec:sim-analysis,sec:sa-program}).

\begin{cppcode}[A projected slab as a PNG]
#include <FML/Visualization/Visualization.h>
using namespace FML::VISUALIZATION;
Image image = project_grid(density, 2, 0.45, 0.1);   // a slab of 10% of the box along z (mean)
if (FML::ThisTask == 0)
    write_png("density.png", image, "magma", ColorScale::automatic_log());
\end{cppcode}

\section{Other utilities}\label{sec:io-other}

\file{ExampleData} generates the small inputs the examples and tests use (a clustered particle box, linear spectra and
tables) in code, so that no data files are needed.

\code{Matplotlibcpp} (adapted from matplotlib-cpp) makes matplotlib plots from C++. \cls{Chunkpix} (by Albert Izard) stores sparse HEALPix
\cite{Gorski2005} maps: the $12N_{\rm side}^2$ pixels are split into $12N_{\rm side,chunk}^2$ chunks and only the chunks that are written to are
allocated; with NEST ordering a chunk is a compact patch, so a partial-sky map allocates about that fraction. It holds the lightcone maps of
\file{SimulationCode} (\cref{ch:weak-lensing}).

\begin{validation}[Validation: I/O and utilities]
\file{tests/test_gadget.cpp}: write and read back on several tasks (positions, peculiar velocities, IDs, header), several species and their
masses, 8-byte IDs, small and empty files, per-particle masses, big- and little-endian files and format 2 with an extra block, HDF5 (including
SWIFT's three-element box size and snapshots with different numbers of files). \file{tests/test_ramses.cpp}: synthetic RAMSES snapshots
(domain reading, units for both Mpc conventions of RAMSES, empty CPU files, a particle at $x=1$, the descriptor file). \file{tests/test_lua.cpp}:
many unset parameters, defaults, arrays and the wrong-type errors. \file{tests/test_utils.cpp}: \code{loadtxt} edge cases, the random generators,
the memory log with many allocations and threads, timings, \cls{ParameterMap} types and the unit systems.
\file{Visualization/tests}: the projections against direct sums (slabs and pencils, every axis, periodic slabs, 1, 2 and
4 tasks), the particle weights and mass conservation, the colour maps and scales, and the PNG and \code{.npy} files read
back with Python's standard library (every pixel compared).
\end{validation}
